\section{A fully materialized equation from unordered proofs}
\label{sec:direct-materialized}

This section closes a different gap from the low-variable record search.
The height-certificate source previously ended at an unelaborated compiler
interface. We now give an ordinary polynomial with every coefficient and
monomial materialized. It is independent of the K215 transition table and
of Jones's universal polynomial. Its size is large: the purpose is a
complete, reproducible direct baseline, not a claim of a new variable record.

The recommended artifact is the literal ordinary polynomial $Q$ in
\texttt{direct/fixed\_base/equation\_integer.json}. It has 27,916 unrestricted
integer variables, exact degree four, 758,275 nonzero monomials, and
maximum absolute coefficient 320,307,372,009,918. In ZFC,
\[
 \operatorname{Con}(\mathrm{ZFC})\quad\longleftrightarrow\quad
 \neg\exists\boldsymbol{x}\in\mathbb Z^{27916}\;Q(\boldsymbol{x})=0.
\]
The natural-domain companion has 16,036 nonnegative variables. It is
specified exactly as the sum of the squared residuals in
\texttt{fixed\_base/residuals\_affine.json}; rebuilding also emits its
collected sparse polynomial \texttt{equation\_affine.json}. Here
$\mathbb N=\{0,1,2,\ldots\}$. Each sparse entry
\texttt{[coefficient,[indices]]} denotes that coefficient times the product
of the listed variables; repeated indices denote repeated factors and
indices start at zero. All coefficients are literal integers, not bounded
search recipes. The compact archive retains the full integer polynomial
and exact residual/circuit data; rebuilding restores redundant expanded
intermediates and checks their recorded SHA256 hashes.

The derivation is mathematical; we do not claim a newly kernel-checked
end-to-end translation. We first explain the two-stream source and then
replace its height stream by a fixed binary bound. The final measurements
above include this reduction and affine-wire elimination.

\subsection{The fixed proof source}

The inherited file \texttt{fixed\_base/horn\_shared\_source.json} specifies twelve
padded single-conclusion rules. Their local relation is
\[
 H(p,q,c;\boldsymbol w)=\prod_{j=0}^{11}
             \left(\sum_{r\in\mathcal R_j}r(p,q,c,\boldsymbol w)^2\right)=0.
\]
There are seventeen shared nonnegative witnesses. The residual lists
$\mathcal R_j$, including the fixed bundled axiom, are included literally.
Their expanded sums of squares are independently reproduced exactly by
the checker. The historical JSON field name
\texttt{quadratic\_residuals} is not a degree assertion: the bundled branch
also contains quartic residuals. The total product has degree fifty.
The first rule is truth $(0,0,0)$; zero premises pad axioms and one-premise
rules. Target code one is falsity.

For clarity, the calibration does not pass through execution of K215.
The repository theorem \texttt{c222\_iff\_zf\_refutation} in
\texttt{KEndpoint.lean}, followed by
\texttt{zf\_refutation\_iff\_zfc\_refutation} in
\texttt{ConstructibleChoice.lean}, identifies derivability of target one
in this numerical calculus with a syntactic ZFC refutation. The exact
numeric constructors and twelve rules are in \texttt{KCode222.lean} and
\texttt{Code254.lean}. These are dependencies of the mathematical
calibration, not claims that the new Python compiler has been checked by Lean.

A finite family of valid rules has a \emph{flow certificate} for one when,
for every positive label $j$,
\[
 \#\{i:c_i=j\}=[j=1]+\#\{i:p_i=j\}+\#\{i:q_i=j\}.
 \tag{F}
\]
Derivability implies (F) by unfolding a proof tree. Conversely, if one is
not derivable, sum these balances over the nonderivable labels in the
finite support. Each valid single-conclusion rule contributes at most
zero: a nonderivable conclusion has a nonderivable premise. The right
side requires a positive net contribution at one, a contradiction.
This argument permits cycles and repeated premises but uses both
nonnegative multiplicities and the single-conclusion hypothesis.

\subsection{The coefficient-height certificate}

Given $N$ rule instances, put
\[
 A(X)=\prod_{i<N}(1+c_iX),\quad
 B(X)=(1+X)\prod_{i<N}(1+p_iX)(1+q_iX),\quad
 R_N=2\prod_{i<N}((1+p_i)(1+q_i)+c_i).
\]
Every coefficient of $A-B$ has absolute value at most $R_N$.
Consequently, for $M=R_N+1$, equality $A(M)=B(M)$ forces $A=B$:
the least nonzero coefficient of a nonzero polynomial vanishing at $M$
would be a nonzero multiple of $M$ of absolute value smaller than $M$.
The distinct roots $-1/j$ recover (F) from the polynomial identity.

This yields the two-accumulator source
\begin{align*}
&S_0=M+1,\quad S_N=1,\quad R_0=2,\quad M=R_N+1,\\
&S_{i+1}(1+c_iM)=S_i(1+p_iM)(1+q_iM),\\
&R_{i+1}=R_i((1+p_i)(1+q_i)+c_i),\qquad H(p_i,q_i,c_i;\boldsymbol w_i)=0.
\end{align*}
All factors are positive. Telescoping gives $A(M)=B(M)$; the second
accumulator supplies the coefficient bound. Conversely, expand a proof
tree in open-obligation order, compute the complete $R_N$ first, choose
$M=R_N+1$, and let $S_i$ be the product of the open-obligation factors.
There is no circular choice of the coefficient bound.

Encode the two sequences using a common beta scale:
\[
 \beta(x,y,i)=x\bmod(1+(i+1)y).
\]
Every pair of finite natural sequences admits such an encoding with
separate $x,z$ and one common $y$, by choosing $y$ large and divisible by
$N!$, then using the Chinese remainder theorem. The factorial divisibility
here is used only in the existence proof; it is not an additional
universally quantified condition in the final source.

The five external parameters are $N,x,z,y,M_0$, with $M=M_0+1$.
Four beta endpoint equations are imposed. In the one bounded row block,
the four beta values are replaced by four triples $(r,q,s)$ satisfying
\[
 x=qd+r,\qquad d=r+s+1,\qquad d=1+(i+1)y,
\]
with the appropriate code and index. Thus the row has exactly
$20+4\cdot3=32$ natural witnesses: twenty for labels and Horn witnesses,
and twelve for the beta graphs. There is a fixed signed polynomial
$P=e-f$ with $e,f$ having nonnegative coefficients such that the source is
\[
 \text{four endpoint equations}\quad\land\quad
 \forall i<N\;\exists\boldsymbol w\in\mathbb N^{32}\quad
 e(i,\boldsymbol w)=f(i,\boldsymbol w).
 \tag{S}
\]
For this two-stream comparison, $e,f$ can be represented as nonnegative
arithmetic DAGs. The delivered \texttt{fixed\_base/row\_source.json} instead
contains the twenty-nine-witness specialization described below. Both use
the following identities without expanding the large Horn product:
$(e-f)+(g-h)=(e+g)-(f+h)$ and
$(e-f)(g-h)=(eg+fh)-(eh+fg)$.

\subsection{One uniform witness bound and one CRT block}

This materialization uses a small improvement to the generic bounded
universal compiler: one common bound $U$ replaces thirty-two independent
coordinate bounds, and the factorial $(U+1)!$ is computed only once.
Set
\[
 T=(e(N,U,\ldots,U)+f(N,U,\ldots,U)+N)!,\qquad
 D_j=1+(j+1)T,\qquad B=\prod_{j<N}D_j.
\]
The bounded block in (S) is equivalent to the following finite
existential conditions, with encoded coordinates $a_1,\ldots,a_{32}$
and one encoded index $v$:
\begin{align*}
 B&\mid (U+1)!\binom{a_i}{U+1} &&(1\leq i\leq32),\\
 B&\mid 1+(v+1)T,\\
 e(v,\boldsymbol a)&\equiv f(v,\boldsymbol a)\pmod B.
\end{align*}

For soundness, fix $j<N$ and take a prime $p\mid D_j$. Since
$\gcd(p,T)=1$ and $T$ is the displayed factorial, $p$ exceeds its factorial
argument. Divisibility of each falling factorial supplies a residue
$w_i\leq U$ with $a_i\equiv w_i\pmod p$. This statement also handles
truncated natural subtraction: if a vanishing factor arises from
$a_i<t$, choose $w_i=a_i$. The index congruence gives $v\equiv j\pmod p$.
Hence $e(j,\boldsymbol w)\equiv f(j,\boldsymbol w)\pmod p$.
Monotonicity of nonnegative-coefficient polynomials bounds both values
strictly below $p$, so they are equal. For completeness choose $U$ above
all finitely many row witnesses. The $D_j$ are pairwise coprime because
$T$ is divisible by all positive differences below $N$; CRT supplies the
encoded columns and index. The argument is also valid for $N=0$.
The original source's endpoints themselves exclude $N=0$.

This is a uniform-bound specialization of the pinned theorem
\texttt{boundedForallCertificate}. It removes no semantic requirement and
needs no unproved new representation theorem.

\subsection{Eliminating every nonpolynomial operation}

The file \texttt{fixed\_base/build\_direct.py} implements the following explicit
identities. Remainders and quotients always have positive denominators
and are encoded by $n=qd+r$, $d=r+s+1$. Congruence modulo $d$, including
$d=0$, is encoded by $a+du=b+dv$ with two natural quotients. If
$Q=((N+1)T+1)^N+2$, then
\[
 B=\bigl(T^N N!\binom{Q+N}{N}\bigr)\bmod(TQ-1).
\]
The denominator is positive since $T\geq1$ and $Q\geq3$.
The binomial and factorial identities are
\begin{align*}
 \binom nk&=\left\lfloor\frac{(2^n+2)^n}{(2^n+1)^k}\right\rfloor
                \bmod(2^n+1),\\
 r&=(n+2)^{n+2},\quad C=\binom{r+n}{n},\qquad
 n!=\left\lfloor\frac{(r+1)^n+C-1}{C}\right\rfloor.
\end{align*}
They include all boundary cases, including $k>n$ and $n=0$.
The predecessor $C-1$ is represented by $C=C_-+1$; here $C>0$.
The pinned \texttt{FactorialRepresentation.lean} proves these identities,
and \texttt{BoundedForallRepresentation.lean} supplies the product formula.

Natural exponentiation is replaced by its fully explicit finite Pell
certificate, not by an oracle. Write
\[
 \mathcal N(a,x,y):\quad x^2+y^2-a^2y^2=1.
\]
The graph of $(x,y)=(X_n(a),Y_n(a))$, for $a>1$, is given by $n\leq y$
and the disjunction of $x=1,y=0$ with the existence of $u,w,s,t,b$ satisfying
\begin{gather*}
 \mathcal N(a,x,y),\quad \mathcal N(a,u,w),\quad \mathcal N(b,s,t),
 \quad b>1,\quad w>0,\\
 b\equiv1\pmod{4y},\quad b\equiv a\pmod u,\quad y^2\mid w,
 \quad s\equiv x\pmod u,\quad t\equiv n\pmod{4y}.
\end{gather*}
This is exactly the pinned finite Matiyasevich certificate in
\texttt{PellRepresentation.lean}. To represent $m=n^k$, use the disjunction
\[
 (k=0\land m=1)\ \lor\
 \left(k>0\land\left[(n=0\land m=0)\ \lor\
       (n>0\land\mathcal C(m,n,k))\right]\right),
\]
where $\mathcal C$ asserts existence of $w,a,t,z,x,y$ with the displayed
Pell-coordinate graph at $(a,k,x,y)$ and
\begin{gather*}
 x\equiv y(a\mathbin{\dot-}n)+m\pmod t,\qquad
 2an=t+n^2+1,\qquad m<t,\\
 n\leq w,\qquad k\leq w,\qquad \mathcal N(w+1,a,wz).
\end{gather*}
Here $\dot-$ is truncated natural subtraction. Its graph is encoded by
$(a=d+n)$ or $(a\leq n\land d=0)$. These formulas are the exact
\texttt{PowerRepresentation.lean} criterion. Strict and weak inequalities
use natural slack variables; divisibility uses a natural quotient.
Conjunctions are sums of nonnegative residual polynomials, and
disjunctions are their products. Consequently inactive branch witnesses
are harmless. In particular, no positivity assumption about an arbitrary
power base is hidden in the generated polynomial.

The delivered fixed-base call ledger contains 108 power certificates,
108 indexed Pell certificates, 33 binomial calls, three factorial calls,
and 72 division/remainder calls. The ledger records actual input and output
node indices and existential ranges for each call.

\subsection{Why the final polynomial is quartic}

All signed expressions are maintained as differences of natural
arithmetic circuits. Every addition or multiplication node receives a
fresh natural wire and a defining residual of degree at most two.
The final equality residuals are linear in the appropriate output wires.
The delivered polynomial is the sum of the squares of all these residuals.
A zero therefore enforces every circuit gate and every arithmetic
certificate. Conversely, any certificate uniquely supplies all gate
wires. Thus flattening preserves existence of zeros in both directions.
A nonzero squared multiplication residual ensures exact degree four;
the checker also verifies the degree of the fully collected polynomial.

Rebuilding emits \texttt{residuals.json} and \texttt{equation.json}, containing,
respectively, the complete residual list and its fully collected sum of
squares. Independent integer arithmetic reconstructs every gate and every
coefficient. The test suite also expands all twelve Horn branches exactly,
checks row evaluations against a separately written source expression,
checks all zero-base and zero-exponent boundaries, and includes a genuine
nonzero-exponent Pell witness for $1^1=1$. These finite checks support
artifact fidelity; the all-input equivalence rests on the proofs above.

\subsection{The recommended smaller source: a fixed binary height bound}

The completed compiler makes a further source improvement measurable.
Choose an independent maximum label $m$, and impose
$p_i,q_i,c_i\leq m$ by three natural slack variables per row. Since
$m+1\leq2^m$ for every natural $m$,
\[
 A(1)\leq(m+1)^N\leq2^{mN},\qquad
 B(1)\leq2(m+1)^{2N}\leq2^{2mN+1}.
\]
It is therefore enough to set
\[
 M_0=2^{2mN+1},\qquad M=M_0+1.
 \tag{BH}
\]
The entire $R$ stream and its beta code disappear. There is now only the
$S$ stream, with endpoints $\beta(x,y,0)=M_0+2$ and
$\beta(x,y,N)=1$. Its local transition is unchanged. The row has
$20+6+3=29$ witnesses instead of thirty-two: twenty for Horn/labels,
six for two beta triples, and three label upper-bound slacks. One global
power certificate imposes (BH). Thus eight power certificates disappear
net, reducing 116 power calls to 108.

Completeness chooses a finite proof tree, then $N$ and a bound $m$ on
its labels, then $M_0$ by (BH), then the $S$ sequence, and finally its
beta encoding and the uniform CRT witness bound $U$. In particular,
$m$ is \emph{not} identified with $U$: the latter also bounds local
remainders and quotients depending on $M$. Such an identification would
introduce a circular, potentially impossible bound. Soundness uses the
same coefficient-height lemma and flow theorem as before. For $N=0$,
(BH) gives $M_0=2$ and the two endpoints are contradictory.

The folder \texttt{direct/fixed\_base/} gives this source and its complete
polynomials. A second, slightly tighter but three-wire-larger version
uses $M_0=2(m+1)^{2N}$; it is an intermediate comparison, not needed
in the recommended release and not a distinct arithmetic theorem.

\subsection{Eliminating affine wires and allowing unrestricted integers}

Addition gates do not need existential output wires. Substitute their
affine expressions into the next multiplication residual. Scalar products
are affine too and are substituted likewise. Retain a wire only for a
multiplication whose expanded expression is genuinely quadratic in
previous retained wires and semantic inputs. Every residual still has
degree at most two. The substitution is triangular and uniquely
restores every omitted wire. All substituted affine forms have
nonnegative coefficients, so natural witnesses are preserved.

Applied to the fixed-binary-height source, this gives the recommended
natural polynomial (expanded by rebuilding as
\texttt{fixed\_base/equation\_affine.json}):
\begin{center}
\begin{tabular}{lr}
Natural variables & 16,036\\
Exact degree & 4\\
Monomials & 700,865\\
Semantic inputs & 3,960\\
Retained nonlinear wires & 12,076\\
Squared residuals & 12,366\\
Maximum absolute coefficient & 320,307,372,009,917
\end{tabular}
\end{center}
The number of monomials increases: fewer variables is not fewer expanded
terms. The full affine substitution map is rebuilt and checked.

For a polynomial over unrestricted integers, only the 3,960 semantic
inputs need explicit nonnegativity constraints. For each such input $x$,
append
\[
 x-a^2-b^2-c(c+1)=0.
\]
Every natural $x$ has this form: apply the three-square theorem to
$4x+1$, whose representation has one odd and two even coordinates, then
divide the identity by four after removing the odd square's constant.
Conversely $a^2+b^2+c(c+1)\geq0$ for integer $a,b,c$.
Internal circuit wires may remain unrestricted because their triangular
multiplication definitions force the unique nonnegative values after
semantic inputs are nonnegative. Adding the squares of these 3,960
quadratic residuals gives the literal integer-domain polynomial in
\texttt{fixed\_base/equation\_integer.json}:
\[
 \boxed{\text{27,916 unrestricted integer variables, exact degree four,
 758,275 monomials.}}
\]
Its maximum absolute coefficient is 320,307,372,009,918.
Both lack-of-solutions statements are provably equivalent in ZFC to
$\operatorname{Con}(\mathrm{ZFC})$. Neither count competes with the
low-variable branch; the advance is an explicit, independently auditable
ordinary equation from the direct proof source.

\subsection{Scope and limitations}

The direct equation is much larger than the low-variable Pell-based
constructions. It should not be described as a numerical record.
Its contribution is completion: no bounded quantifier, variable exponent,
factorial, binomial coefficient, unknown compiler output, or unspecified
coefficient remains in the final polynomial file. The fixed numeric proof
source, the general arithmetic representations, and the flattening theorem
are separate, auditable layers. Future source optimizations can now be
measured against an actual ordinary equation rather than an interface.
